\section{Morphology-Conditioned Motion Imitation}\label{sec:method}

\subsection{Task, simulation, and actuation}

For a clip $c$ and body configuration $b$, let
$\hat{s}_{c,b}(t)$ be the refined reference state and $s_t$ the simulated state
of the matching asset. We learn a shared policy
$\pi_\theta(a_t\mid o_t,b)$ that tracks this reference on flat ground.
The body assigned to an environment is fixed, while its reference motion can
change on reset. No language input or external scene object is supplied.

Experiments use the \texttt{smpl\_mor} assets in Isaac Gym through
ProtoMotions~\cite{isaacgym,protomotions}. Simulation runs at 60\,Hz with two
substeps and control decimation two, giving a 30\,Hz policy rate. The policy
outputs a raw action $a_t\in\R^D$, where $D=69$. Joint targets are
\begin{equation}
    q_t^{\mathrm{target}}
      = q^{\mathrm{offset}}
        + q^{\mathrm{scale}}\odot\tanh(a_t),
    \label{eq:action_mapping}
\end{equation}
with offsets and scales derived from the configured joint ranges. These
targets drive the simulator's built-in position PD controller. The action is
not a residual added to the reference pose.

PD stiffness, damping, and effort limits are scaled per joint by the mass of
its associated rigid body relative to that body in the reference asset.
This is a body-level mass scaling, not a single scale based on total character
mass, and it does not guarantee dynamically equivalent actuation across
geometrically different bodies. Base gains and limits are listed in
Table~\ref{tab:pd_parameters}.

\subsection{Observation design}\label{sec:observations}

The observation combines present state, history, future reference targets,
a historical action, and explicit morphology information.
Let $B=24$ denote the number of simulated bodies. The state representation
contains root height above ground, root-relative body positions excluding
the root, 6D body rotations, and linear/angular velocities.
Positions and velocities use a heading-aligned frame, and rotations remove
the root heading. This removes dependence on arbitrary global heading and
horizontal origin while retaining height and body configuration.

We include seven historical states at control-step offsets
$\{1,2,3,4,8,16,32\}$, covering approximately 1.07\,s, and four future
reference slots at $\{1,2,4,8\}$, covering approximately 0.27\,s.
Each historical state uses its own root/heading frame.
Future slots contain target positions relative to the current root,
heading-aligned target rotations, position differences from the current body
state, relative rotations $\smash{R_{t,j}^{-1}\hat{R}_{t+k,j}}$, and
heading-aligned linear/angular velocity differences. Thus relative rotations
are expressed with respect to the current body rotation, not as another
heading-frame position feature.

\begin{table}[t]
\centering
\small
\caption{Observation inputs supplied to both actor and critic. Dimensions are
per slot. Positions/heights use meters, linear velocities m/s, angular
velocities rad/s, and 6D rotations are dimensionless. Raw actions and
morphology coefficients are not joint-angle measurements.}\label{tab:observations}
\begin{tabularx}{\linewidth}{@{}lrlY@{}}
\toprule
Input & Dimension & Slots & Representation and path \\
\midrule
Current state & $15B-2=358$ & 1 &
Height; $3(B-1)$ positions; $6B$ rotations; $3B$ linear and $3B$ angular
velocities. Current-state encoder. \\
History & $358$ & 7 &
Same state representation at the seven past offsets. Shared history encoder. \\
Future reference & $24B=576$ & 4 &
Target and relative poses ($18B$), plus velocity differences ($6B$).
Shared future encoder. \\
Historical action & $D=69$ & 1 &
A raw previous policy-action vector, before tanh and PD scaling.
Bypasses attention. \\
Morphology & $11$ & 1 &
$[g,\beta_1,\ldots,\beta_{10}]$ from the assigned asset.
Bypasses attention. \\
\bottomrule
\end{tabularx}
\end{table}

The temporal context exposes both recent changes and more distant motion
information; the future targets provide a short anticipation window.
These are design motivations, not established ablation outcomes.
History is initialized consistently with reference-state resets using
available historical reference states; single-state initialization is the
fallback when a trajectory history is unavailable.
Contact flags, an explicit phase scalar, text, and estimated physical
feature vectors are not part of the selected observation.

\subsection{Temporal slot/type attention}\label{sec:architecture}

We use three observation encoders: one for the current state, one shared
across history slots, and one shared across future slots. Each is an MLP
with two 256-unit ReLU hidden layers and a 256-dimensional token output.
Inputs use running mean/standard-deviation normalization and a magnitude
clamp of five.

The resulting 12 tokens comprise one current, seven historical, and four
future tokens. For slot $i$, the Transformer input is
\begin{equation}
    z_i = E_{\tau(i)}(x_i)
        + e_i^{\mathrm{slot}} + e_{\tau(i)}^{\mathrm{type}},
    \qquad i=0,\ldots,11,
    \label{eq:token_embedding}
\end{equation}
where $\tau(i)$ identifies the current/history/future source.
Slot embeddings distinguish the configured temporal offsets, while type
embeddings distinguish the observation sources. The Transformer has two
layers, four attention heads, token width 256, and feed-forward width
1,024~\cite{attention}. Its current-state token supplies the pooled output
after the configured output projection and activation.

\begin{figure}[t]
\centering
\fbox{\begin{minipage}{0.94\linewidth}
\small
\begin{tabularx}{\linewidth}{@{}YcY@{}}
Current state: $1\times358$ & $\longrightarrow$ &
Current encoder: $1\times256$ token \\
Past states: $7\times358$ & $\longrightarrow$ &
History encoder: $7\times256$ tokens \\
Future targets: $4\times576$ & $\longrightarrow$ &
Future encoder: $4\times256$ tokens
\end{tabularx}
\begin{center}
$\downarrow$\\
12 tokens + learned temporal-slot and source-type embeddings\\
$\downarrow$\\
2 Transformer layers, 4 heads, width 256\\
$\downarrow$\\
Pool current token: $h_t\in\R^{256}$\\
$\downarrow$\\
Concatenate $[h_t,\ a^{\mathrm{hist}}_t,\ g,\boldsymbol{\beta}]$
and normalize
\end{center}
\begin{tabularx}{\linewidth}{@{}YY@{}}
\textbf{Actor:} six 2,896-unit layers
$\longrightarrow$ 69 action means
$\longrightarrow$ tanh/range scaling
$\longrightarrow$ PD targets. &
\textbf{Critic:} four 1,024-unit layers
$\longrightarrow$ scalar state value.
\end{tabularx}
\smallskip
The entire encoder/attention pathway is instantiated separately for actor
and critic; the drawing summarizes the common design, not shared weights.
Morphology and action bypass attention into each head. Text is absent.
\end{minipage}}
\caption{Observation-to-action data flow of the selected slot/type policy.
Both networks receive the same input categories but have separate parameters.}\label{fig:network}
\end{figure}

The final head input is the concatenation of the 256-dimensional attention
output, the historical raw action, and the 11-dimensional morphology vector.
This 336-dimensional input also uses running normalization and clamp five.
Betas are not manually divided by three, embedded as a separate attention
token, or used for adaLN modulation in the selected model.

Actor and critic have separate token encoders, attention modules, slot/type
embeddings, and MLP heads. The actor head has six hidden layers of width
2,896; the critic has four of width 1,024, all with ReLU activations.
The outputs are a Gaussian action mean and a scalar value, respectively.
PPO samples actions with fixed diagonal log standard deviation $-2.9$;
evaluation uses the mean action. Figure~\ref{fig:network} summarizes the flow.

\subsection{Tracking reward, resets, and PPO}\label{sec:learning}

The training reward combines global body tracking and root-height tracking.
For corresponding simulated and reference body positions $p_j,\hat{p}_j$,
rotations $R_j,\hat{R}_j$, and velocities $v_j,\hat{v}_j$ and
$\omega_j,\hat{\omega}_j$, define
\begin{align}
    e_p &= \frac{1}{3B}\sum_{j=1}^{B}\|p_j-\hat{p}_j\|_2^2,
    & e_R &= \frac{1}{B}\sum_{j=1}^{B}
                  d_{\SO(3)}(R_j,\hat{R}_j)^2, \\
    e_v &= \frac{1}{3B}\sum_{j=1}^{B}\|v_j-\hat{v}_j\|_2^2,
    & e_\omega &= \frac{1}{3B}\sum_{j=1}^{B}
                     \|\omega_j-\hat{\omega}_j\|_2^2, \\
    e_h &= (h-\hat{h})^2, &&
\end{align}
where $d_{\SO(3)}$ is the geodesic rotation angle in radians.
The coordinate-wise means in the position and velocity errors match the
implemented reductions. The reward is
\begin{equation}
\begin{split}
    r_t ={}& 0.5\exp(-25e_p) + 0.3\exp(-5e_R)
             + 0.1\exp(-0.5e_v)\\
           &+0.2\exp(-0.1e_\omega) + 0.2\exp(-100e_h).
\end{split}\label{eq:reward}
\end{equation}
The weights are not normalized to sum to one. There is no explicit effort,
action-smoothness, or contact-matching reward in this experiment, and there
is no source-dependent switch between canonical and HUMOS reward functions.

References are sampled only from variants matching the environment's asset.
Reference-state initialization starts at frame zero with probability 0.2;
otherwise it uses a random valid reference time. Training episodes have a
maximum of 1,000 control steps and use the configured fall termination with
a 0.15\,m height setting and permitted foot-contact bodies. Tracking-error
termination is disabled. This training rule is distinct from the evaluation
failure threshold in Section~\ref{sec:metrics}.

We optimize using PPO~\cite{ppo}, with 32-step rollouts, discount 0.99,
GAE parameter 0.95, policy-ratio clipping 0.2, and normalized advantages.
Separate Adam optimizers use learning rates $2\times10^{-5}$ for the actor
and $10^{-4}$ for the critic. Other configuration values and actual run
budgets appear in Table~\ref{tab:training_parameters}. Evaluation rollouts
use deterministic inference without gradient tracking.

\subsection{Direct multi-shape learning and streamed sampling}\label{sec:sampling}

The final approach is direct multi-shape training on refined HUMOS
references. A small static corpus supports initial development; the
full-scale experiment streams per-clip files through \texttt{GlobalClipPool}.
The training manifest is deterministically partitioned across GPU ranks.
Each rank maintains a bounded resident pool containing all 128 variants
of its selected clips. Environment-level motion sampling respects the
assigned morphology.

At each scheduled rebuild, most resident slots are selected by explicit
top priority, combining estimated clip difficulty and an exploration bonus.
A fraction 0.2 is reserved for uniform random selection, maintaining rehearsal
outside the highest-priority set. Difficulty estimates are updated from
training-pool evaluations with an exponential moving average of rate 0.1
and a floor of 0.05. The exploration coefficient is 1.0.
This is a training curriculum, not uniform sampling of the complete corpus.

The resident pool is rebuilt every 256 epochs. The full-scale run initially
used 256 clips per rank and was resumed with 128 after memory failures; the
local cache multiplier was reduced to 1.0. These changes affect residency
and memory demand, not the underlying training split. Validation membership
is fixed independently of pool rebuilding, and validation evaluation does
not update training sampling weights. We report resident-pool and validation
metrics separately because they measure different populations.
